% PRISM method section draft for a CoRL-style paper.
% Required packages in the main paper preamble:
% \usepackage{amsmath,amssymb,booktabs,tikz}
% \usetikzlibrary{arrows.meta,positioning}

\section{Method}
\label{sec:method}

We propose PRISM, a Prefix-Routed Indexed Slot Memory module for long-horizon
imitation learning. PRISM is designed as a plug-and-play history module rather
than as a policy-specific architectural branch. Its role is to maintain a
fixed-budget streaming summary of the executed prefix and to expose this summary
through a conditioning interface that can be consumed by an action expert, such
as an ACT-style chunking Transformer, a diffusion policy, or an autoregressive
controller. The central idea is to use a state-prefix signature as an explicit
addressing signal for a slot-structured visual memory: the signature does not
only provide an additional context token, but directly routes where and how the
visual prefix is written.

\subsection{Problem Formulation}
\label{sec:method_problem}

We assume access to a demonstration dataset
\begin{equation}
    \mathcal{D} = \{\tau_i\}_{i=1}^{N}, \qquad
    \tau_i = \{(o_t^i, s_t^i, a_t^i)\}_{t=1}^{T_i},
\end{equation}
where $o_t$ denotes the multi-view visual observation, $s_t$ denotes the
low-dimensional robot state, and $a_t \in \mathbb{R}^{d_a}$ is the expert action. Let
$x_t=(o_t,s_t)$ be the current observation. For an action-chunking expert with
horizon $H$, the supervised target is
\begin{equation}
    y_t^\star = (a_t^\star, a_{t+1}^\star,\ldots,a_{t+H-1}^\star).
\end{equation}
A standard action expert predicts $\hat y_t=\mathcal{E}_{\psi}(x_t)$ from the current
observation alone or from a short observation window. In contrast, PRISM
augments the expert with a learned prefix condition:
\begin{equation}
    \hat y_t =
    \mathcal{E}_{\psi}\!\left(x_t;\, c_t\right),
    \qquad
    c_t = \mathcal{C}_{\theta}(g_t,\Delta g_t,M_t,z_t^{\mathrm{mem}}).
    \label{eq:prism_policy}
\end{equation}
Here $g_t$ is a projected state-prefix signature, $\Delta g_t$ is its projected
increment, $M_t \in \mathbb{R}^{K \times d}$ is a fixed-size slot memory, and
$z_t^{\mathrm{mem}}$ is an optional pooled memory readout. The action expert
keeps its native prediction head and training objective; PRISM supplies only
the additional history condition $c_t$.

\begin{figure*}[t]
    \centering
    \begin{tikzpicture}[
        font=\scriptsize,
        >=Latex,
        node distance=0.72cm and 0.88cm,
        block/.style={draw, rounded corners=2pt, align=center, minimum height=0.72cm, minimum width=1.55cm},
        wide/.style={draw, rounded corners=2pt, align=center, minimum height=0.72cm, minimum width=2.25cm}
    ]
        \node[block] (state) {state prefix\\$s_{1:t}$};
        \node[block, below=of state] (obs) {visual prefix\\$o_{1:t}$};
        \node[wide, right=of state] (sig) {truncated signature\\$\xi_t,\delta_t$};
        \node[wide, right=of obs] (enc) {visual/state encoders\\$v_t,p_t$};
        \node[wide, right=of sig, yshift=-0.72cm] (router) {prefix router\\$r_t,\alpha_t$};
        \node[wide, right=of router] (slots) {indexed slot memory\\$M_t=[m_t^{(1)},\ldots,m_t^{(K)}]$};
        \node[wide, right=of slots] (adapter) {context adapter\\$c_t$};
        \node[wide, right=of adapter] (expert) {action expert\\$\mathcal{E}_{\psi}$};
        \node[block, right=of expert] (action) {action chunk\\$\hat y_t$};

        \draw[->] (state) -- (sig);
        \draw[->] (sig) -- (router);
        \draw[->] (obs) -- (enc);
        \draw[->] (state.east) to[bend right=20] (enc.west);
        \draw[->] (enc) -- (router);
        \draw[->] (router) -- (slots);
        \draw[->] (slots) -- (adapter);
        \draw[->] (sig.east) to[bend left=15] (adapter.north west);
        \draw[->] (enc.east) to[bend right=15] (adapter.south west);
        \draw[->] (adapter) -- (expert);
        \draw[->] (expert) -- (action);
        \draw[->, dashed] (slots.south) .. controls +(0,-0.72) and +(0,-0.72) .. node[below] {streaming state} (router.south);
    \end{tikzpicture}
    \caption{
    PRISM overview. A state-prefix signature produces a routing signal that
    indexes a fixed number of visual memory slots. The resulting slot memory,
    signature tokens, and pooled readout are converted into an expert-specific
    context while leaving the action expert architecture otherwise unchanged.
    }
    \label{fig:prism_overview}
\end{figure*}

\subsection{State-Prefix Signature}
\label{sec:method_signature}

Let $S_{1:t}$ denote the continuous interpolation of the state trajectory
$(s_1,\ldots,s_t)$. We encode this prefix with a truncated path signature of
depth $p$:
\begin{equation}
    \xi_t =
    \operatorname{Sig}_{\leq p}(S_{1:t})
    =
    \left(
    \int_{1<u_1<\cdots<u_m<t}
    dS_{u_1}\otimes\cdots\otimes dS_{u_m}
    \right)_{m=1}^{p}
    \in \mathbb{R}^{d_{\mathrm{sig}}}.
    \label{eq:path_signature}
\end{equation}
We also use the first-order change in the prefix descriptor,
\begin{equation}
    \delta_t = \xi_t - \xi_{t-1}, \qquad \xi_0=\mathbf{0}.
    \label{eq:delta_signature}
\end{equation}
Both quantities are projected to the policy hidden dimension,
\begin{equation}
    g_t = \phi_g(\xi_t), \qquad
    \Delta g_t = \phi_{\Delta}(\delta_t),
    \qquad g_t,\Delta g_t \in \mathbb{R}^d.
    \label{eq:signature_projection}
\end{equation}
During training, $\xi_t$ and $\delta_t$ are computed from demonstration prefixes.
During deployment, they are updated online from the executed state trajectory, so
the descriptor at time $t$ contains no future information.

\subsection{Prefix-Routed Indexed Slot Memory}
\label{sec:method_prism_memory}

\paragraph{Visual and state encoders.}
At each step, the current visual observation and state are embedded as
\begin{equation}
    v_t = \rho\!\left(f_{\mathrm{vis}}(o_t)\right), \qquad
    p_t = \phi_s(s_t),
    \label{eq:vis_state_embed}
\end{equation}
where $f_{\mathrm{vis}}$ is the visual backbone, $\rho$ pools spatial and camera
features, and $p_t,v_t \in \mathbb{R}^d$. PRISM maintains $K$ memory slots,
\begin{equation}
    M_t =
    [m_t^{(1)},\ldots,m_t^{(K)}] \in \mathbb{R}^{K \times d},
    \qquad M_0=\mathbf{0}.
\end{equation}

\paragraph{Signature-indexed routing.}
The prefix signature is first converted to a routing state:
\begin{equation}
    r_t = \phi_r([g_t,\Delta g_t]) \in \mathbb{R}^{d_r},
    \label{eq:route_state}
\end{equation}
where $\Delta g_t$ is omitted when delta routing is disabled. The routing query
and slot keys are
\begin{equation}
    q_t = W_q r_t, \qquad
    k_{t-1}^{(j)} = W_k m_{t-1}^{(j)}.
\end{equation}
The write distribution over memory slots is then
\begin{equation}
    \ell_t^{(j)}
    =
    \frac{q_t^\top k_{t-1}^{(j)}}{\sqrt{d_r}},
    \qquad
    \alpha_t^{(j)}
    =
    \frac{\exp(\ell_t^{(j)})}{\sum_{l=1}^{K}\exp(\ell_t^{(l)})}.
    \label{eq:slot_routing}
\end{equation}
This routing step is the key distinction of PRISM: the slot address is a
function of the executed state path, not only of the current image.

\paragraph{Slot-wise gated writing.}
We form a shared write proposal
\begin{equation}
    u_t = \phi_w([v_t,p_t,g_t,\Delta g_t]) \in \mathbb{R}^d,
    \label{eq:write_base}
\end{equation}
and compute per-slot candidates and gates:
\begin{align}
    \tilde m_t^{(j)}
    &=
    \phi_m([m_{t-1}^{(j)},u_t,r_t]),
    \\
    \gamma_t^{(j)}
    &=
    \sigma\!\left(\phi_{\gamma}([m_{t-1}^{(j)},u_t,r_t])\right).
\end{align}
For a valid prefix element, the slot update is
\begin{equation}
    m_t^{(j)}
    =
    m_{t-1}^{(j)}
    +
    \alpha_t^{(j)}\gamma_t^{(j)}
    \left(\tilde m_t^{(j)} - m_{t-1}^{(j)}\right).
    \label{eq:slot_update}
\end{equation}
For padded prefix elements during training, the update is masked and
$m_t^{(j)}=m_{t-1}^{(j)}$. Equation~\eqref{eq:slot_update} separates routing,
write strength, and content transformation, making the memory state
interpretable as a set of prefix-indexed caches rather than a single mixed
recurrent hidden state.

\paragraph{Memory readout.}
For expert adapters that require a pooled context, PRISM reads the slot memory
with an attention query conditioned on the current observation and route state:
\begin{align}
    q_t^{\mathrm{read}}
    &=
    \phi_q^{\mathrm{read}}([v_t,p_t,r_t]),
    \\
    \beta_t^{(j)}
    &=
    \operatorname{softmax}_j
    \left(
    \frac{(q_t^{\mathrm{read}})^\top W_k^{\mathrm{read}}m_t^{(j)}}{\sqrt{d}}
    \right),
    \\
    z_t^{\mathrm{mem}}
    &=
    \sum_{j=1}^{K}
    \beta_t^{(j)} W_v^{\mathrm{read}}m_t^{(j)}.
    \label{eq:memory_readout}
\end{align}
When this readout is not used, the pooled context is replaced by the mean of the
slots. The full PRISM state is therefore fixed-size, requiring
$O(Kd+d_{\mathrm{sig}})$ storage independent of the horizon length.

\subsection{Action Expert Adapter}
\label{sec:method_adapter}

PRISM exposes the generic context in Eq.~\eqref{eq:prism_policy}; the adapter
maps this context into the format expected by the action expert. In our ACT
instantiation, the expert predicts an action chunk with Transformer decoder
queries. The current observation is encoded into Transformer memory
$H_t^{\mathrm{enc}}$, while the PRISM context is injected as extra decoder
cross-attention memory tokens:
\begin{equation}
    \mathcal{K}_t =
    [\, e_t^{g}, e_t^{\Delta}, m_t^{(1)},\ldots,m_t^{(K)}, H_t^{\mathrm{enc}} \,],
    \label{eq:act_memory_tokens}
\end{equation}
where $e_t^{g}=g_t$ and $e_t^{\Delta}=\Delta g_t$ after projection. If the model
uses memory-conditioned encoder FiLM, the pooled context also modulates
current-observation encoder tokens:
\begin{equation}
    [\Gamma_t,B_t] = \Psi_{\mathrm{FiLM}}(z_t^{\mathrm{mem}}),
    \qquad
    \tilde X_t =
    X_t \odot (1+\tanh \Gamma_t) + B_t.
    \label{eq:film}
\end{equation}
The ACT decoder then predicts
\begin{equation}
    \hat y_t =
    \operatorname{Dec}_{\psi}(Q_H,\mathcal{K}_t),
    \qquad
    \hat y_t \in \mathbb{R}^{H \times d_a},
\end{equation}
where $Q_H$ are the learned chunk queries. Other experts can consume the same
PRISM state through different adapters, for example as diffusion conditioning
embeddings, autoregressive prefix tokens, cross-attention keys and values, or
RNN hidden-state initializers.

\begin{table}[t]
    \centering
    \caption{Examples of PRISM adapters for common imitation-learning experts.}
    \label{tab:prism_adapters}
    \begin{tabular}{lll}
        \toprule
        Expert & Adapter & Condition form \\
        \midrule
        ACT / chunking Transformer & extra memory tokens, FiLM & tokens, modulation \\
        Diffusion policy & denoising conditioner & vector or context tokens \\
        Autoregressive policy & prefix or KV adapter & prefix tokens, KV memory \\
        MLP / RNN BC & concatenation or initialization & vector, hidden state \\
        \bottomrule
    \end{tabular}
\end{table}

\subsection{Training}
\label{sec:method_training}

For a training target at time $t$, we construct an explicit prefix subsequence
$\tilde P_t=(q_1,\ldots,q_L)$ from $\{1,\ldots,t\}$ under a fixed budget $L$ and
stride $r$. The construction always keeps the current step as the last valid
prefix element; when the prefix exceeds the budget, it preserves the first
episode step and the most recent tail. The sequence is right-padded and
accompanied by a Boolean mask $b_l$ indicating whether $q_l$ is valid.

Training reconstructs the same memory state that would have been obtained
online by scanning the prefix in temporal order:
\begin{equation}
    M_{q_l}
    =
    \mathcal{U}_{\theta}(M_{q_{l-1}},o_{q_l},s_{q_l},g_{q_l},\Delta g_{q_l}; b_l),
    \qquad l=1,\ldots,L.
    \label{eq:prefix_scan}
\end{equation}
The final valid state is used to condition the action expert at time $t$. This
scan prevents future leakage: padded elements do not update the memory, and the
last valid prefix element is exactly the supervised current step.

For ACT, the primary imitation objective is an action-chunk regression loss with
an action-padding mask $\Omega$:
\begin{equation}
    \mathcal{L}_{\mathrm{BC}}
    =
    \frac{1}{|\Omega|}
    \sum_{(t,h)\in\Omega}
    \left\|
    \hat a_{t,h} - a_{t+h}^{\star}
    \right\|_1.
    \label{eq:bc_loss}
\end{equation}
When the VAE-style ACT objective is enabled, we add the standard KL term:
\begin{equation}
    \mathcal{L}_{\mathrm{expert}}
    =
    \mathcal{L}_{\mathrm{BC}}
    +
    \lambda_{\mathrm{KL}}
    D_{\mathrm{KL}}
    \left(
    q_{\psi}(z\mid s_t,y_t^\star)\,\|\,\mathcal{N}(0,I)
    \right).
    \label{eq:expert_loss}
\end{equation}
PRISM can also use optional memory-specific regularizers. A load-balancing
loss discourages all writes from collapsing into one slot:
\begin{equation}
    \bar \alpha^{(j)}
    =
    \frac{1}{N_{\mathrm{valid}}}
    \sum_{t,l} b_l \alpha_{q_l}^{(j)},
    \qquad
    \mathcal{L}_{\mathrm{bal}}
    =
    \sum_{j=1}^{K}
    \left(
    \bar \alpha^{(j)} - \frac{1}{K}
    \right)^2.
    \label{eq:balance_loss}
\end{equation}
A read-write consistency loss aligns the pooled memory readout with the current
write proposal:
\begin{equation}
    \mathcal{L}_{\mathrm{cons}}
    =
    \frac{1}{N_{\mathrm{valid}}}
    \sum_{t,l} b_l
    \left\|
    z_{q_l}^{\mathrm{mem}} - u_{q_l}
    \right\|_2^2.
    \label{eq:consistency_loss}
\end{equation}
The final objective is
\begin{equation}
    \mathcal{L}
    =
    \mathcal{L}_{\mathrm{expert}}
    +
    \lambda_{\mathrm{bal}}\mathcal{L}_{\mathrm{bal}}
    +
    \lambda_{\mathrm{cons}}\mathcal{L}_{\mathrm{cons}}.
    \label{eq:total_loss}
\end{equation}
The auxiliary coefficients may be set to zero, in which case PRISM is trained
only through the native imitation-learning loss of the action expert.

\begin{figure*}[t]
    \centering
    \begin{tikzpicture}[
        font=\scriptsize,
        >=Latex,
        node distance=0.7cm and 0.9cm,
        block/.style={draw, rounded corners=2pt, align=center, minimum height=0.72cm, minimum width=1.75cm},
        wide/.style={draw, rounded corners=2pt, align=center, minimum height=0.72cm, minimum width=2.35cm}
    ]
        \node[block] (train_prefix) {training prefix\\$\tilde P_t$};
        \node[wide, right=of train_prefix] (train_scan) {masked temporal scan\\Eq.~\eqref{eq:prefix_scan}};
        \node[wide, right=of train_scan] (train_context) {PRISM context\\$c_t$};
        \node[wide, right=of train_context] (train_loss) {expert loss\\$+\;$optional memory losses};

        \node[block, below=1.15cm of train_prefix] (test_obs) {online step\\$x_t$};
        \node[wide, right=of test_obs] (test_update) {one-step update\\$M_{t-1}\mapsto M_t$};
        \node[wide, right=of test_update] (test_context) {PRISM context\\$c_t$};
        \node[wide, right=of test_context] (test_action) {execute action\\from $\hat y_t$};

        \draw[->] (train_prefix) -- (train_scan);
        \draw[->] (train_scan) -- (train_context);
        \draw[->] (train_context) -- (train_loss);
        \draw[->] (test_obs) -- (test_update);
        \draw[->] (test_update) -- (test_context);
        \draw[->] (test_context) -- (test_action);
        \draw[->, dashed] (train_scan.south) -- node[right] {same update rule} (test_update.north);
        \node[align=left, above=0.08cm of train_prefix] {Training};
        \node[align=left, above=0.08cm of test_obs] {Inference};
    \end{tikzpicture}
    \caption{
    Training and inference consistency. Training reconstructs the fixed-budget
    memory by scanning an explicit, masked prefix. Inference maintains the same
    state online with a one-step PRISM update after each observation.
    }
    \label{fig:prism_train_infer}
\end{figure*}

\subsection{Online Inference}
\label{sec:method_inference}

At the beginning of each episode, PRISM resets $M_0=\mathbf{0}$ and clears the
signature history. At step $t$, the policy appends the current state to the
history, computes $\xi_t$ and $\delta_t$, updates $M_t$ with
Eq.~\eqref{eq:slot_update}, and calls the action expert with the resulting
context $c_t$. For chunking experts, the policy may execute the first
$n_{\mathrm{act}}$ actions from the predicted chunk; setting
$n_{\mathrm{act}}=1$ gives the most responsive use of the updated streaming
memory. Since PRISM stores only $K$ slots and the current signature descriptor,
its memory footprint does not grow with episode length.
